\Article{4}
The problem set by the Royal Society of Sciences requires the image to be similar to the object mapped in the smallest parts. The first step is to express this condition analytically.

Let differentiation of the functions of \(t,u\) that express \(x,\allowbreak y,\allowbreak z,\allowbreak X,\allowbreak Y,\allowbreak Z\) give the following equations:
\[
\begin{aligned}
\mathrm{d}x&=a\,\mathrm{d}t+a'\,\mathrm{d}u\\
\mathrm{d}y&=b\,\mathrm{d}t+b'\,\mathrm{d}u\\
\mathrm{d}z&=c\,\mathrm{d}t+c'\,\mathrm{d}u\\
\mathrm{d}X&=A\,\mathrm{d}t+A'\,\mathrm{d}u\\
\mathrm{d}Y&=B\,\mathrm{d}t+B'\,\mathrm{d}u\\
\mathrm{d}Z&=C\,\mathrm{d}t+C'\,\mathrm{d}u
\end{aligned}
\]
\noindent The prescribed condition requires, first, that all infinitesimal lines issuing from one point on the first surface and lying in it be proportional to their corresponding lines on the second surface, and second, that the former make the same angles with one another as the latter.

Such a line element on the first surface is
\[
\begin{aligned}
={}&\BareRoot{\bigl((aa+bb+cc)\mathrm{d}t^2+2(aa'+bb'+cc')\mathrm{d}t\mathord{.}\mathrm{d}u}\\
&\qquad+(a'a'+b'b'+c'c')\mathrm{d}u^2\bigr)
\end{aligned}
\]
\noindent and the corresponding element on the second surface is
\[
\begin{aligned}
={}&\BareRoot{\bigl((AA+BB+CC)\mathrm{d}t^2+2(AA'+BB'+CC')\mathrm{d}t\mathord{.}\mathrm{d}u}\\
&\qquad+(A'A'+B'B'+C'C')\mathrm{d}u^2\bigr).
\end{aligned}
\]
\noindent If the two are to stand in a definite ratio to one another, independently of \(\mathrm{d}t\) and \(\mathrm{d}u\), then the three quantities
\[
aa+bb+cc,\quad aa'+bb'+cc',\quad a'a'+b'b'+c'c'
\]
\noindent must evidently be proportional, respectively, to the following three:
\SourceLeaf{200}
\[
AA+BB+CC,\quad AA'+BB'+CC',\quad A'A'+B'B'+C'C'
\]
\noindent If the values
\[
t,u\quad\text{and}\quad t+\delta t,u+\delta u
\]
\noindent correspond to the endpoints of a second element on the first surface, then the cosine of the angle it makes with the first element is
\[
=\frac{\begin{gathered}(a\,\mathrm{d}t+a'\,\mathrm{d}u)(a\delta t+a'\delta u)+(b\,\mathrm{d}t+b'\,\mathrm{d}u)(b\delta t+b'\delta u)\\+(c\,\mathrm{d}t+c'\,\mathrm{d}u)(c\delta t+c'\delta u)\end{gathered}}{\begin{gathered}\BareRoot{((a\,\mathrm{d}t+a'\,\mathrm{d}u)^2+(b\,\mathrm{d}t+b'\,\mathrm{d}u)^2+(c\,\mathrm{d}t+c'\,\mathrm{d}u)^2)}\\\mathord{.}((a\delta t+a'\delta u)^2+(b\delta t+b'\delta u)^2+(c\delta t+c'\delta u)^2)\end{gathered}}
\]
\noindent and an entirely similar expression for the cosine of the angle between the corresponding elements on the second surface is obtained simply by changing \(a,b,c,a',b',c'\) into \(A,B,C,A',B',C'\). Evidently the two expressions become equal when the above proportionality holds, so the second condition is already included in the first; this also becomes clear of itself on some reflection.

The analytical expression of the condition of our problem is therefore that
\[
\frac{AA+BB+CC}{aa+bb+cc}=\frac{AA'+BB'+CC'}{aa'+bb'+cc'}=\frac{A'A'+B'B'+C'C'}{a'a'+b'b'+c'c'}
\]
\noindent must hold. This will be a finite function of \(t\) and \(u\), which we shall set \(=mm\). Then \(m\) expresses the ratio by which linear dimensions on the first surface are enlarged or reduced in their image on the second (according as \(m\) is greater or less than 1). Generally speaking, this ratio will vary from place to place: in the particular case where \(m\) is constant, there will be perfect similarity also in the finite parts, and if, in addition, \(m=1\), there will be perfect equality, and one surface will be developable on the other.

